
\subsection*{5.5. Complements}
\phantomsection
\label{sga3:VIIB:5.5}
\addcontentsline{toc}{subsection}{5.5. Complements}
\footnote{\label{sga3:VIIB:nde:169}The editors added this
subsection in order to give some consequences of \SGARefLink{sga3:vii:corollary:5-4}{5.4} mentioned
in Expos\'es III and VIA.}

Recall the following definitions. A Noetherian local ring \(A\)
is called a \emph{complete intersection} if its completion
\(\widehat A\) is the quotient of a complete regular Noetherian
local ring \(B\) by an ideal \(I\) generated by a regular
sequence of elements of \(B\) (cf.\ EGA IV\(_4\), 19.3.1).

On the other hand, let
\[
  \tau:Y\hookrightarrow X
\]
be a closed immersion of schemes. If \(y\in Y\), then \(\tau\)
is called a \emph{regular immersion at the point \(y\)} if the
kernel of
\[
  \mathcal O_{X,y}\longrightarrow\mathcal O_{Y,y}
\]
is generated by a regular sequence. If, moreover, \(X\) is
locally Noetherian and \(\tau\) is a regular immersion at every
point, then \(\tau\) is called a \emph{regular immersion}
(cf.\ \textit{loc.\ cit.}, Proposition~16.9.10 and
Definition~16.9.2).

\medskip
\noindent\textbf{Corollary 5.5.1.}
\phantomsection
\label{sga3:VIIB:5.5.1}
\textit{If \(G\) is an algebraic group over a field \(k\), then
the local ring \(\mathcal O_{G,e}\) is a complete intersection.}

\medskip
\noindent\textit{Proof.}
By EGA IV\(_4\), 19.3.4, we may assume that \(k\) is
algebraically closed. If \(\operatorname{char}(k)=0\), then
\(G\) is already known to be smooth (cf.\ \SGARefLink{sga3:VIIA:3.3.1}{3.3.1} or VIB, \SGARefLink{cor:sga3-VIB-1.6.1}{1.6.1});
hence \(\widehat{\mathcal O}_{G,e}\) is a \(k\)-algebra of
formal power series, by EGA IV\(_4\), 17.5.3
\((\mathrm d^{\prime\prime})\). If
\(\operatorname{char}(k)=p>0\), then \SGARefLink{sga3:vii:corollary:5-4}{5.4}, applied with
\(H=\{e\}\), shows that \(\widehat{\mathcal O}_{G,e}\) is a
complete intersection.

\medskip
\noindent\textbf{Remarks 5.5.2.}
\phantomsection
\label{sga3:VIIB:5.5.2}
\textit{Let \(k\) be a field, let \(G\) be a smooth algebraic
\(k\)-group, and let \(H\) be a closed subgroup of \(G\).}

\smallskip
\noindent\textit{(a)}
We saw in \SGARefLink{sga3:iii:proposition:4-15}{Expos\'e III, 4.15}, that the immersion
\[
  H\hookrightarrow G
\]
is regular. This can also be deduced from \SGARefLink{sga3:vii:corollary:5-4}{5.4} as follows. As in
\textit{loc.\ cit.}, we may assume that \(k\) is algebraically
closed, and it is enough to show that the kernel \(I\) of
\[
  \mathcal O_{G,e}\longrightarrow\mathcal O_{H,e}
\]
is generated by a regular sequence. Put
\[
  A=\mathcal O_{G,e},
  \qquad
  \widehat G=\operatorname{Spf}(\widehat A),
  \qquad
  \widehat H
  =
  \operatorname{Spf}(\widehat{\mathcal O}_{H,e}).
\]
Since \(A\) is Noetherian, there is an exact sequence
\[
  0
  \longrightarrow I\otimes_A\widehat A
  \longrightarrow \widehat A
  \xrightarrow{\ \pi\ }
  \mathscr A(\widehat H)
  \longrightarrow 0,
\]
and \(\widehat A\) is faithfully flat over \(A\). Therefore, by
EGA IV\(_4\), 16.9.10(ii) and 19.1.5(ii), it is enough to show
that the kernel
\[
  \widehat I=I\otimes_A\widehat A
\]
of \(\pi\) is generated by a regular sequence of elements of
\(\widehat A\).

Since \(G\) is smooth, \(\widehat A\) is reduced. By \SGARefLink{sga3:vii:corollary:5-4}{5.4}, the
subalgebra
\[
  B=\mathscr A(\widehat G/\widehat H)
\]
is therefore isomorphic to a formal power series algebra
\[
  k[[x_1,\ldots,x_n]].
\]
Consequently, the identity section of
\(\widehat G/\widehat H\) is defined in \(B\) by the regular
sequence \((x_1,\ldots,x_n)\). Since \(\widehat A\) is
Noetherian, the ideal \(J\) of \(\widehat A\) generated by
\(x_1,\ldots,x_n\) is closed, and hence equals \(\widehat I\),
by Corollary~1.4. Moreover, \(\widehat A\) is topologically
flat, and hence flat over \(B\) (cf.\ \SGARefLink{sga3:VIIB:0.3.8}{0.3.8}). Thus
\((x_1,\ldots,x_n)\) is a regular sequence in \(\widehat A\),
by EGA IV\(_4\), 19.1.5(ii).

\smallskip
\noindent\textit{(b)}
The following more precise result can also be deduced from
\SGARefLink{sga3:VIIB:5.2:theorem:item:ii}{5.2(ii)}. Suppose that \(k\) is perfect. Applying \SGARefLink{sga3:VIIB:5.2:theorem:item:ii}{5.2(ii)} to the
algebra
\[
  C=\mathscr A(H),
\]
there are a basis
\[
  (y_1,\ldots,y_{r+s})
\]
of \(\omega_H\) and integers
\[
  1\leq n_1\leq\cdots\leq n_r
\]
such that \(\mathscr A(H)\) is isomorphic to the quotient of
\[
  k[[y_1,\ldots,y_{r+s}]]
\]
by the ideal generated by
\[
  y_i^{p^{n_i}},
  \qquad i=1,\ldots,r.
\]
Lift the \(y_i\) to elements \(x_i\) of \(\omega_G\), and extend
\((x_1,\ldots,x_{r+s})\) to a basis
\((x_1,\ldots,x_n)\) of \(\omega_G\). Since
\(\mathscr A(G)\) is reduced, the morphism
\[
  k[[x_1,\ldots,x_n]]
  \longrightarrow
  \mathscr A(G)
\]
is an isomorphism, by \SGARefLink{sga3:VIIB:5.2:theorem:item:iii}{5.2(iii)}.

Thus there exists a ``coordinate system''
\((x_1,\ldots,x_n)\) on \(G\), that is, an isomorphism
\[
  \mathscr A(G)
  \simeq
  k[[x_1,\ldots,x_n]],
\]
such that \(H\) is defined by the equations
\[
  x_i^{p^{n_i}}=0
  \quad (i=1,\ldots,r),
  \qquad
  x_i=0
  \quad (i>r+s).
\]
This is ``Dieudonn\'e's theorem''; compare \textup{[Di55]},
\(\S\)19, Theorem~6, and \textup{[Di73]}, II, \(\S\)\SGARefLink{sga3:VIIB:3.2:prop}{3.2},
Proposition~3 and the material preceding it.

